\section{Extensions and open items}\label{sec:app:extensions}

The yield-ceiling variant that previously headed this list --- the soft ceiling
$\lim_{x\to\infty}a\left(x,t\right)=1$ --- is no longer an extension: it is carried in the model as
case~(S) of Section~\ref{sec:sp:setup:primitives}, alongside the minimum material requirement
$\bar R\ge0$. What remains here are variants we have not built.

\smalltitle{Residence time beyond the waste stock.} The waste stock $\mathcal W$ of
Section~\ref{sec:sp:setup:primitives} gives the recovery pipeline its residence time and ties
circular throughput to the stock of matter actually in the economy. Three refinements remain
unbuilt. First, residence time in \emph{use}: the embodied content of consumption, of extraction and
exploration effort, and of treatment is still released as waste in the instant the final good is
absorbed; giving every use $j$ a stock $M^j$ with release rate $\delta^j$ would replace the pair
$\left(\Phi^j,\Omega^j\right)$ by a single intensity plus a release rate, with the current
specification recovered as $\delta^j\to\infty$, and the ledger survives unchanged in form with
$M^K$ replaced by total material in use. Second, the handling rate $\mu$ as a control --- an
explicit landfill-mining margin, faster recovery of the stockpile at rising marginal cost --- which
makes $\mathcal W$ a reserve depleted at a chosen rate rather than a passive buffer. Third, a
leaking stockpile: a term $\ell\mathcal W$ in $\dot P$ that gives landfills a flow cost alongside
their asset value. The theory keeps $\mu$ constant and the stockpile inert; the quantitative model
is the natural place for the last two margins.

\smalltitle{A scale-dependent material requirement.} The minimum material requirement $\bar R$ is a
constant flow, which is the strong form of the assumption: it says the floor is independent of how
large the economy is. The natural alternative is a maintenance reading, $\bar R=\bar r K$, under which
the floor is the material needed to hold the capital stock intact --- a close relative of the
$\Phi^K\delta K$ term the accounting already carries. This is endogenous and shrinks as capital
shrinks, so it softens the finite-duration result of Section~\ref{sec:sp:longrun} into a trade-off
between the size of the capital stock and the length of the material era, and it is the version a
reader is most likely to ask about.

\smalltitle{Irreversible pollution.} Assumption~\ref{ass:lr:theta} imposes a regeneration floor,
$\theta\left(P\right)\ge\theta_{\min}>0$, making pollution fully reversible: once emissions cease the
stock regenerates completely. Permanent legacies --- $\theta\to0$ at finite $P$ --- would add a
hysteresis dimension to every entry of the long-run classification, and interact with the
non-convexity the decay function already introduces.

\smalltitle{Alternative embodied-materials assumptions.} The two variants of $\Phi^I$ set aside in
Section~\ref{sec:sp:setup:aggregatematerial} --- exogenous in levels, and constant over time --- are
recorded in Appendices~\ref{sec:app:exogenousPhiI} and~\ref{sec:app:constantPhiI}.

\smalltitle{One resource stock.} Collapsing extraction and exploration into a single activity is the
special case $C^D_{DD}=0$ (just-in-time discovery). It is not the model we want --- the wedge between
the reserve rent and the discovery margin is the economics of the resource side, and the
stranding-versus-exhaustion distinction is unstatable with one stock --- but it is a streamlined
presentation worth having available.
